\section*{Cluster A Step 24: Robustness Stress of Group-Shape Generation}

\paragraph{Grade.}
Finite GUT-exterior closure audit / recognition recovery stress test. The result is not a physical gauge-sector theorem, and the positive Step-23 landing is demoted because the exterior-role ansatz and slot-five condition remain inputs.

\paragraph{Generated-vs-input.}
Inputs are the widened factor/rank/U(1)-role range, the known-alternative stress set, closure criteria, minimality variations, the declared dual/antisymmetric/mixed exterior-role ansatz, and the resulting \( \mathrm{total\_slots}=5 \) condition. Outputs are the closure table, stability map, and fragile verdict within that declared grammar.

\paragraph{Wider closure table.}
\[
\begin{array}{c|ccc}
n & \mathrm{candidates} & \mathrm{closers} & \mathrm{closing\ structures}\\
\hline
0 & 4 & 0 & \varnothing\\
1 & 27 & 3 & SU(5),SO(10),SU(6)\\
2 & 84 & 1 & 2|3\\
3 & 224 & 0 & \varnothing\\
4 & 504 & 0 & \varnothing\\
5 & 1008 & 0 & \varnothing\\
6 & 1848 & 0 & \varnothing
\end{array}
\]

\paragraph{Criteria variation.}
The split-sector \(2|3\) winner is stable in \(8/9\) criteria settings. The factor-count-first ordering selects \(SU(5)\), so uniqueness is criteria-dependent.

\paragraph{Known alternatives.}
\(SU(5)\), \(SO(10)\), and the \(SU(6)\) chiral example are inserted recognition controls with assigned content-pattern labels and anomaly status. They are not generated by the Step-24 base closure engine.

\paragraph{Verdict.}
The verdict is \(\mathrm{FRAGILE}\): Step 23 is hardened against the original smuggle but not against all reasonable criteria. The next grammar must distinguish low-rank split-sector closure from one-factor simple unification economy.
